\documentclass{report}
\usepackage{amsmath,amssymb}
\setlength{\parindent}{0mm}
\setlength{\parskip}{1em}
\begin{document}
\begin{center}
\rule{6in}{1pt} \
{\large Arvind K. Saibaba \\
{\bf Fast algorithms for the Geostatistical approach for solving Linear Inverse problems based on Hierarchical matrices}}

Institute for Computational &Mathematical Engg \\ Jen-Hsun Huang Engineering Center \\ #53T \\ 475 Via Ortega \\ Stanford \\ CA 94305-4121
\\
{\tt arvindks@stanford.edu}\\
Peter K. Kitanidis\end{center}

In several inverse problems that arise from geophysical applications,
such as identifying the contaminant source from time history of spatially
distributed contaminant measurements, the prevalent approach is to use
the Geostatistical approach with Gaussian priors. There are two main
computational bottlenecks in the large-scale implementation of this
approach: (1) covariance matrices that arise from finely spaced
discretizations on irregular grids, can be extremely large and moreover,
dense. (2) for certain problems with a large number of measurements, the
measurement operator is not only dense, forming it explicitly would
require the repeated solution of (possibly) time-dependent partial
differential equations. In this talk we propose to show, how to deal with
each one of these issues.\\

The covariance matrices that arise in practice, are extremely large and
dense and this usually places a heavy burden both in terms of storage and
computational requirements. Covariance kernels which are stationary and
translation invariant, discretized on a uniformly spaced, regular grid,
result in Toeplitz (or Block-Toeplitz) matrices. Toeplitz matrices can be
embedded inside circulant matrices, for which operations such as
matrix-vector products, matrix-matrix products can be efficiently
computed using Fast Fourier Transforms (FFT). However, their primary
deficiency is that these algorithms don't extend very easily to other
types of grids that are predominant in realistic problems. Covariance
matrices, although dense, have special structure which can be exploited.
They are similar to dense matrices that arise from the discretization of
integral equations. Using the Hierarchical matrix approach, we will show
how to reduce the storage and computational complexity of matrix-vector
products to ${\cal{O}}(N\log N)$, where $N$ is the number of unknowns
that we are solving for. \\

Then, using Bayes' rule, we can write down the posterior probability
distribution of the unknowns, given the measurements, and assuming a
Gaussian prior for the unknowns. The maximum a posteriori (MAP) estimate
can be computed by solving the system of equations. We propose to use a
matrix-free Krylov subspace approach to solve the resulting system of
equations. This approach has a huge computation advantage because it
avoids the explicit construction of the matrix and only relies on
matrix-vector products, which can be accelerated using the Hierarchical
matrix approach. We also propose a preconditioner that serves to cluster
the eigenvalues and therefore reduce the number of iterations taken by
the iterative solver. We will also provide numerical evidence for the
clustering of the eigenvalues. Another important aspect of the
Geostatistical approach is that, it not only allows us to obtain the best
estimate via maximizing the likelihood, we can also quantify the
uncertainty associated with our estimate. This can be done by generating
conditional realizations from our posterior probability distribution of
the unknowns. We show how to do this efficiently by using Chebyshev
matrix polynomial approximation for the square root of the covariance
matrix, to compute unconditional realizations and using these
unconditional realizations to generate conditional realizations by
solving a modification of the system of equations used to compute the MAP
estimate. Finally, we demonstrate the performance of our algorithm on a
large-scale inverse model problem on unstructured grids from contaminant
source identification.


\end{document}
